% paper5_global_regularity_route2.tex
% T-First Programme — Paper 5
% "Global Regularity for the Route 2 Navier–Stokes–Theta System at eps > 0"
% Target: Annals of Mathematics
% Status: FRAMEWORK DRAFT — analytical gaps identified, proof strategy set
%         Depends on: Paper 3 (Claim A in 2D), M7 (3D numerical evidence)
%         Primary open step: Estimate (4.8) — ε-suppression of stretching
%
\documentclass[11pt]{article}
\usepackage{amsmath,amssymb,amsthm}
\usepackage{geometry}
\usepackage{hyperref}
\usepackage{microtype}
\usepackage{booktabs}
\usepackage{xcolor}
\usepackage{enumitem}
\geometry{margin=1.2in}

% ── Theorem environments ───────────────────────────────────────────────────────
\newtheorem{theorem}{Theorem}[section]
\newtheorem{proposition}[theorem]{Proposition}
\newtheorem{lemma}[theorem]{Lemma}
\newtheorem{corollary}[theorem]{Corollary}
\theoremstyle{definition}
\newtheorem{definition}[theorem]{Definition}
\newtheorem{assumption}[theorem]{Assumption}
\newtheorem{remark}[theorem]{Remark}
\theoremstyle{remark}
\newtheorem{claim}[theorem]{Claim}
\newtheorem{openproblem}[theorem]{Open Problem}

% ── Macros ─────────────────────────────────────────────────────────────────────
\newcommand{\R}{\mathbb{R}}
\newcommand{\T}{\mathbb{T}}
\newcommand{\dx}{\,dx}
\newcommand{\dt}{\,dt}
\newcommand{\norm}[1]{\left\|#1\right\|}
\newcommand{\abs}[1]{\left|#1\right|}
\newcommand{\Lp}[1]{L^{#1}}
\newcommand{\Hs}[1]{H^{#1}}
\newcommand{\QT}{Q_T}
\newcommand{\prt}{\partial_t}
\newcommand{\grad}{\nabla}
\newcommand{\divg}{\operatorname{div}}
\newcommand{\curl}{\operatorname{curl}}
\newcommand{\mueff}{\mu_{\mathrm{eff}}}
\newcommand{\epsst}{\varepsilon^*}
\newcommand{\epsp}{\varepsilon_{\mathrm{param}}}
\newcommand{\dlt}{\delta}
\newcommand{\Sth}{S_\theta}
\newcommand{\Feps}{\mathcal{F}_\eps}

% ── Title block ────────────────────────────────────────────────────────────────
\title{%
  Global Regularity for the Route~2\\
  Navier--Stokes--Theta System at $\epsp>0$\\[6pt]
  {\normalsize\color{red}\textbf{[FRAMEWORK DRAFT — T-First Programme, Paper 5]}}\\
  {\normalsize\color{red}Primary gap: Estimate~\eqref{eq:key-estimate} (Section~\ref{sec:gap})}
}

\author{%
  Pratik Menon\thanks{Email: pmenonn@gmail.com}
  \\[4pt]
  \small T-First Programme
}

\date{May 2026}

% ══════════════════════════════════════════════════════════════════════════════
\begin{document}
\maketitle

\begin{abstract}
We present the framework for a global regularity theorem for the Route~2
system of the T-First programme:
\begin{align*}
  \prt u + (u\cdot\grad)u &= -\grad p + \nu\Delta u + \epsp\divg(f(\theta)\grad u),
  \quad \divg u = 0,\\
  \prt\theta + u\cdot\grad\theta &= \nu\Delta\theta + \nu|\grad u|^2,
  \quad \theta(\cdot,0) = 0,
\end{align*}
on $\T^3\times[0,T]$ with $\nu>0$ and $\epsp>0$.  The system reduces to the
exact Prize Navier--Stokes equations at $\epsp=0$.

\medskip

The main theorem we aim to prove is:

\begin{center}
\textbf{Theorem P5 (target):} \emph{For every smooth $u_0$ with $\divg u_0=0$,
there exists $\epsst=\epsst(E_0,\nu)>0$ such that for $\epsp>\epsst$, the
Route~2 system has a globally smooth solution on $[0,\infty)$.}
\end{center}

The proof has four phases.  Phases~0 and~1 are complete or near-complete.
Phase~2 is in progress.  Phase~3 depends on the central open estimate
(\S\ref{sec:gap}).  The numerical evidence from M7 ($128^3$) strongly supports
the theorem for $\epsst<0.001$.

\medskip

The primary analytical gap is: controlling the vorticity stretching integral
$\int(\omega\cdot\grad)u\cdot|\omega|^{p-2}\omega\dx$ by the Route~2 extra
dissipation $\epsp\int f(\theta)|\grad u|^2\dx$.  Two candidate routes
(CZ + $\eps$-suppression; spectral multiplier via the $\theta$-equation) are
analysed in detail.
\end{abstract}

\tableofcontents

% ══════════════════════════════════════════════════════════════════════════════
\section{Introduction}
\label{sec:intro}

\subsection{The Target Theorem}

This paper aims to prove:

\begin{theorem}[Global regularity for Route~2, $\epsp>0$]\label{thm:main}
  Let $u_0\in C^\infty(\T^3)$ with $\divg u_0=0$, and let $\nu>0$.
  There exists a threshold $\epsst=\epsst(\norm{u_0}_{H^1},\nu)>0$ such that
  for every $\epsp\geq\epsst$, the Route~2 system \eqref{eq:NS-R2}--\eqref{eq:theta}
  has a unique globally smooth solution $(u,\theta,p)$ on $\T^3\times[0,\infty)$.
\end{theorem}

\begin{remark}
  The threshold $\epsst$ depends on the initial energy $E_0=\frac{1}{2}\norm{u_0}^2_{L^2}$
  and viscosity $\nu$.  It does NOT depend on $T$ — the bound is global in time.
  The numerical experiments (M7, \S\ref{sec:numerics}) suggest $\epsst<0.001$
  for $E_0\sim 0.125$, $\nu=0.001$.
\end{remark}

\subsection{Why This Suffices for the Prize (Paper~6)}

Theorem~\ref{thm:main} combined with the $\epsp\to 0$ limit argument
(Paper~6) gives:
\begin{equation}
  u^\epsp \xrightarrow{\epsp\to 0} u^0, \qquad
  u^0 \text{ solves exact Prize NS and is smooth.}
\end{equation}
The limit is non-singular: at $\epsp=0$ the velocity equation is the exact
Prize equation.  Compactness (Aubin--Lions) gives convergence; lower
semi-continuity of the LPS margin gives smoothness of $u^0$.

\subsection{The Route~2 System}
\label{sec:system}

The Route~2 system on $\T^3\times(0,\infty)$:
\begin{align}
  \prt u + (u\cdot\grad)u &= -\grad p + \nu\Delta u
    + \epsp\,\divg(f(\theta)\,\grad u),
    \quad \divg u = 0, \label{eq:NS-R2}\\[4pt]
  \prt\theta + u\cdot\grad\theta &= \nu\Delta\theta + \nu|\grad u|^2,
    \quad \theta(\cdot,0) = 0, \label{eq:theta}\\[4pt]
  u(\cdot,0) &= u_0 \in C^\infty(\T^3), \quad \divg u_0 = 0. \label{eq:IC}
\end{align}
Here $f:\R\to[0,\infty)$ is smooth, non-decreasing, with $f(0)=0$,
$f'(s)\geq 0$, $f(s)\leq C_f(1+s)$ for all $s\geq 0$ (linear growth).
The standard choice is $f(\theta)=\theta$.

The \emph{effective viscosity} is $\mueff = \nu + \epsp f(\theta)\geq\nu>0$.

\subsection{The Four-Phase Proof Structure}

\begin{table}[h]
\centering
\caption{Four-phase proof structure for Theorem~\ref{thm:main}.}
\begin{tabular}{@{}cllll@{}}
\toprule
Phase & Goal & Key tool & Status & Gap \\
\midrule
0 & Local existence + continuation criterion & Energy + compactness & Complete & None \\
1 & $\theta\in L^\infty(Q_T)$ & $L^{n/2+\eps}$ Claim A + parabolic theory & Near-complete & Claim A in 3D \\
2 & Vorticity $L^p$ bound at $\epsp>0$ & $\eps$-suppression of stretching & In progress & Estimate \eqref{eq:key-estimate} \\
3 & Global strong solutions & Continuation criterion never triggered & Follows from 0+1+2 & Phase 2 \\
\bottomrule
\end{tabular}
\end{table}

% ══════════════════════════════════════════════════════════════════════════════
\section{Phase~0: Local Existence and Continuation}
\label{sec:phase0}

\subsection{Local Existence}

\begin{theorem}[Local existence]\label{thm:local}
  For $u_0\in\Hs{3}(\T^3)$, there exists $T^*=T^*(u_0,\nu,\epsp)>0$ and a
  unique smooth solution $(u,\theta)$ to \eqref{eq:NS-R2}--\eqref{eq:IC}
  on $[0,T^*]$.
\end{theorem}

\begin{proof}[Proof sketch]
  Standard energy method + Picard iteration in $\Hs{3}$.  The $\theta$
  equation is a parabolic equation with $L^2$-source (from energy bound on $u$),
  giving local existence in $L^2_t\Hs{2}_x$.  The system is locally
  Lipschitz in $\Hs{3}\times\Hs{2}$, so a contraction mapping argument closes.
\end{proof}

\subsection{Continuation Criterion}

\begin{theorem}[Blowup criterion]\label{thm:blowup}
  If $T^*<\infty$ (finite-time blowup), then necessarily
  \begin{equation}\label{eq:blowup-criterion}
    \limsup_{t\to T^*} \norm{\omega(\cdot,t)}_{L^3(\T^3)} = \infty,
  \end{equation}
  where $\omega=\curl u$ is the vorticity.  Equivalently, $u\notin L^p_t L^q_x(Q_{T^*})$
  for any Prodi--Serrin pair $(p,q)$ with $2/p+3/q\leq 1$.
\end{theorem}

\begin{proof}
  Standard Beale--Kato--Majda criterion \cite{BKM1984} applies since the
  Route~2 system has $\mueff\geq\nu>0$ (uniformly parabolic).  The $\theta$
  equation is subcritical given the $u$-regularity, so $\theta$ does not
  create additional blowup.
\end{proof}

\begin{remark}
  Theorem~\ref{thm:main} is proved if and only if \eqref{eq:blowup-criterion}
  never holds — i.e., $\norm{\omega}_{L^3}$ remains bounded for all time.
  The global regularity problem is thus equivalent to an a priori bound on
  the $L^3$ vorticity.
\end{remark}

% ══════════════════════════════════════════════════════════════════════════════
\section{Phase~1: $\theta$ is Globally Bounded}
\label{sec:phase1}

\subsection{Non-negativity and $L^1$ Bound}

\begin{lemma}[Positivity of $\theta$]\label{lem:theta-pos}
  Under \eqref{eq:theta}, $\theta(x,t)\geq 0$ for all $(x,t)$.
\end{lemma}

\begin{proof}
  The source $\nu|\grad u|^2\geq 0$ and $\theta(\cdot,0)=0$.  By the
  comparison principle for parabolic equations, $\theta\geq 0$.
\end{proof}

\begin{lemma}[$L^1$ bound on $\theta$]\label{lem:theta-L1}
  \begin{equation}
    \norm{\theta(\cdot,t)}_{L^1(\T^3)}
    \leq \nu\int_0^t\norm{\grad u(\cdot,s)}^2_{L^2}\ds
    \leq E_0/\nu.
  \end{equation}
\end{lemma}

\begin{proof}
  Integrate \eqref{eq:theta} over $\T^3$.  The transport and Laplacian
  terms vanish; the source gives $\frac{d}{dt}\int\theta\dx = \nu\int|\grad u|^2\dx$.
  Integrate in time and use the energy inequality.
\end{proof}

\subsection{$L^\infty$ Bound Conditional on Claim~A}

\begin{proposition}[$\theta\in L^\infty$ from Claim~A]\label{prop:theta-Linfty}
  Suppose Claim~A holds in 3D with $\delta\geq 1/2$, i.e.,
  $\norm{S_\theta}_{L^{3/2}(Q_T)}<\infty$.  Then $\theta\in L^\infty(Q_T)$
  with
  \begin{equation}
    \norm{\theta}_{L^\infty(Q_T)} \leq C(\nu,E_0,T).
  \end{equation}
\end{proposition}

\begin{proof}
  By parabolic $L^q$-theory (Lemma~2.3 of Paper~3) with $q=3/2>1$:
  $\theta\in L^{3/2}_t W^{2,3/2}_x$.  By the Sobolev embedding
  $W^{2,3/2}(\T^3)\hookrightarrow L^\infty(\T^3)$ (since $2\cdot(3/2)>3$),
  $\theta(\cdot,t)\in L^\infty(\T^3)$ for a.e.~$t$.  The bound is uniform
  in $t$ by the parabolic maximal regularity estimate.
\end{proof}

\begin{remark}
  The threshold $\delta\geq 1/2$ (Claim~A with exponent $3/2$) is the
  critical value for Phase~1 in 3D.  The M7 experiments test whether
  $\delta_{\max}\geq 1/2$ at $N=128^3$.  From $N=32$ and $N=64$ runs,
  $\delta_{\max}=1.50\gg 1/2$, strongly suggesting Phase~1 is accessible.
\end{remark}

% ══════════════════════════════════════════════════════════════════════════════
\section{Phase~2: Vorticity $L^p$ Bound — The Central Gap}
\label{sec:gap}

\subsection{The Vorticity Equation in Route~2}

Taking $\curl$ of \eqref{eq:NS-R2}:
\begin{equation}\label{eq:vort-R2}
  \prt\omega + (u\cdot\grad)\omega
  = \nu\Delta\omega + (\omega\cdot\grad)u
  + \epsp\,\curl\bigl(\divg(f(\theta)\grad u)\bigr).
\end{equation}
The last term is the Route~2 \emph{extra dissipation contribution} to vorticity.
We write it as $\mathcal{R}_\epsp[\theta,u]$.

\begin{lemma}[Extra dissipation term]\label{lem:extra-dissip}
  \begin{equation}
    \mathcal{R}_\epsp[\theta,u]
    = \epsp\bigl[f(\theta)\Delta\omega + (\grad f(\theta)\cdot\grad)\omega
      + \curl(\grad f(\theta)\times\grad u)\bigr].
  \end{equation}
  In particular, $\int\mathcal{R}_\epsp\cdot|\omega|^{p-2}\omega\dx$
  contains a negative semi-definite term $-\epsp f(\theta)\int|\grad\omega|^2|\omega|^{p-2}\dx\leq 0$
  (dissipation) plus lower-order terms involving $\grad f(\theta)$.
\end{lemma}

\subsection{The Vorticity $L^p$ Energy Identity}

Multiplying \eqref{eq:vort-R2} by $|\omega|^{p-2}\omega$ and integrating:
\begin{equation}\label{eq:vort-Lp-identity}
  \frac{1}{p}\frac{d}{dt}\norm{\omega}^p_{L^p}
  = -\nu C_p\norm{\omega}^{p-2}_{L^p}\norm{\grad\omega}^2_{L^p}
    - \epsp D_p[\theta,\omega]
    + I_{\mathrm{stretch}}(t)
    + R_p[\theta,\omega],
\end{equation}
where:
\begin{align*}
  D_p[\theta,\omega] &:= C_p\int f(\theta)|\grad\omega|^2|\omega|^{p-2}\dx \geq 0
    \quad\text{(Route~2 extra dissipation)},\\
  I_{\mathrm{stretch}}(t) &:= \int(\omega\cdot\grad)u\cdot|\omega|^{p-2}\omega\dx
    \quad\text{(stretching --- sign unknown)},\\
  R_p[\theta,\omega] &:= \int\curl(\grad f(\theta)\times\grad u)\cdot|\omega|^{p-2}\omega\dx
    \quad\text{(lower-order remainder)}.
\end{align*}

\subsection{The Primary Open Estimate}
\label{sec:open-estimate}

Theorem~\ref{thm:main} requires closing \eqref{eq:vort-Lp-identity} —
showing the right-hand side is non-positive for all time.  The dissipation
terms are negative; the transport term vanishes (incompressibility).
The only obstruction is the stretching integral $I_{\mathrm{stretch}}$.

\begin{openproblem}[Primary gap — $\eps$-suppression of stretching]\label{op:main}
  Prove that for $\epsp>\epsst(E_0,\nu)>0$:
  \begin{equation}\label{eq:key-estimate}
    I_{\mathrm{stretch}}(t)
    \leq \frac{\nu C_p}{2}\norm{\omega}^{p-2}_{L^p}\norm{\grad\omega}^2_{L^p}
      + \epsp\,D_p[\theta,\omega] + C_{\mathrm{abs}}(\nu,E_0,p),
  \end{equation}
  where $C_{\mathrm{abs}}$ is an \emph{absorbing} constant (bounded uniformly in $t$).
  Estimate \eqref{eq:key-estimate} would imply $\norm{\omega}_{L^p}$ is bounded
  for all $t$, closing the Phase~2 vorticity bound.
\end{openproblem}

\begin{remark}
  If \eqref{eq:key-estimate} holds, the Gr\"onwall inequality applied to
  \eqref{eq:vort-Lp-identity} gives
  $\norm{\omega(\cdot,t)}_{L^p}\leq C(\nu,E_0,p,t)$ for all $t>0$,
  and Theorem~\ref{thm:main} follows immediately.
\end{remark}

% ── Section 4.4: Two candidate routes ─────────────────────────────────────────
\subsection{Route~B: CZ + $\eps$-Suppression}
\label{sec:routeB}

\textbf{Step 1 (Bound $I_{\mathrm{stretch}}$ by $\norm{\omega}^{p+1}$).}
By H\"older + Sobolev:
\begin{equation}\label{eq:stretch-Lp}
  \abs{I_{\mathrm{stretch}}} \leq \norm{(\omega\cdot\grad)u}_{L^p}\norm{\omega}^{p-1}_{L^p}
  \leq C\norm{\grad u}_{L^6}\norm{\omega}^{p-1}_{L^p}\norm{\omega}_{L^{6p/(6-p)}}
  \quad (p<6).
\end{equation}
CZ gives $\norm{\grad u}_{L^6}\leq C\norm{\omega}_{L^6}$.
Interpolation $\norm{\omega}_{L^6}\leq C\norm{\omega}^{1-\alpha}_{L^2}\norm{\grad\omega}^\alpha_{L^2}$
($\alpha=3/4$) then gives:
\begin{equation}
  \abs{I_{\mathrm{stretch}}}
  \leq C\norm{\omega}^{p-1}_{L^p}\cdot\norm{\omega}^{1-\alpha}_{L^2}
    \cdot\norm{\grad\omega}^\alpha_{L^2}.
\end{equation}

\textbf{Step 2 (Absorb by dissipation).}
The goal is to absorb the stretching by
$\nu C_p\norm{\omega}^{p-2}_{L^p}\norm{\grad\omega}^2_{L^p}$ (standard dissipation)
or $\epsp D_p[\theta,\omega]$ (Route~2 extra dissipation).

Young's inequality ($ab\leq \eps_Y a^{1/\alpha} + C(\eps_Y)b^{1/(1-\alpha)}$) applied to the
stretching bound gives:
\begin{equation}\label{eq:Young}
  \abs{I_{\mathrm{stretch}}}
  \leq \frac{\nu C_p}{4}\norm{\grad\omega}^2_{L^p}\norm{\omega}^{p-2}_{L^p}
    + C(\nu,E_0,p)\,\norm{\omega}^\gamma_{L^p}
\end{equation}
for some $\gamma<p$ (provided $p>3$, using $\norm{\omega}_{L^2}\leq C$ from energy).

\textbf{Step 3 (Close via absorbing constant).}
The remaining term $C\norm{\omega}^\gamma_{L^p}$ with $\gamma<p$ is subcritical.
By Gr\"onwall + the standard energy bound:
\begin{equation}
  \norm{\omega}^\gamma_{L^p} \leq C_1 + C_2\norm{\omega}^p_{L^p}/(\text{absorbing})
\end{equation}
This closes the loop if the Route~2 extra dissipation $\epsp D_p[\theta,\omega]$
is large enough to absorb $C_2\norm{\omega}^p_{L^p}$.

\textbf{Critical condition (Route~B):}
\begin{equation}\label{eq:routeB-condition}
  \epsp\,D_p[\theta,\omega] \geq C(\nu,E_0,p)\,\norm{\omega}^p_{L^p}
  \quad \Leftrightarrow \quad
  \epsp\,f(\theta_{\min}) \geq C'(\nu,E_0,p).
\end{equation}
Since $f(\theta_{\min})\geq 0$ depends on the dynamics, this requires
a lower bound on $\theta_{\min}$ at positive times — which follows from
Proposition~\ref{prop:theta-Linfty} (Phase~1) if Claim~A holds with $\delta\geq 1/2$.

\begin{remark}[Route~B closure condition]
  Route~B closes if $\epsp>C'/f(\theta_{\min})$.  Since $\theta_{\min}>0$
  for $t>0$ (Lemma~\ref{lem:theta-pos} + strong maximum principle), and
  $f>0$ for $\theta>0$, we have $f(\theta_{\min})>0$.  The condition
  $\epsp>C'/f(\theta_{\min})$ is satisfiable for any fixed $t>0$, but
  we need a uniform-in-$t$ lower bound on $\theta_{\min}$.  This is the
  remaining gap in Route~B.
\end{remark}

\subsection{Route~C: Spectral Multiplier via the $\theta$-Equation}
\label{sec:routeC}

An alternative approach uses the $\theta$-equation directly.  Rearranging
\eqref{eq:theta}:
\begin{equation}\label{eq:Stheta-identity}
  \nu|\grad u|^2 = \Sth = \prt\theta + u\cdot\grad\theta - \nu\Delta\theta.
\end{equation}
Each term on the right has better integrability than $L^1$:
\begin{itemize}
\item $\prt\theta$: from parabolic theory, in $L^q$ for some $q>1$ if the source is.
\item $\nu\Delta\theta$: $\theta\in L^2_t\Hs{1}_x$ from energy → $\Delta\theta\in L^2_tH^{-1}_x$.
  But with Claim~A, $\theta\in L^2_t\Hs{2}_x$, giving $\Delta\theta\in L^2$.
\item $u\cdot\grad\theta$: requires Hölder estimate $\norm{u}_{L^r}\norm{\grad\theta}_{L^s}$.
\end{itemize}

The Route~C argument is:
\begin{enumerate}
\item Start with $\Sth\in L^{1+\delta_0}$ (Claim~A, some $\delta_0>0$).
\item Parabolic regularity: $\theta\in L^{1+\delta_0}_tW^{2,1+\delta_0}_x$.
\item Sobolev: $\theta\in L^{1+\delta_0}_t L^{q(\delta_0)}_x$ for some $q(\delta_0)>6$.
\item Identity \eqref{eq:Stheta-identity}: $\Sth = \prt\theta + u\cdot\grad\theta - \nu\Delta\theta$.
  Bound each term in $L^{1+\delta_1}$ for $\delta_1>\delta_0$.
\item Iterate: $\delta_n\to\infty$.
\end{enumerate}

The critical estimate in Route~C is bounding $u\cdot\grad\theta$:
\begin{equation}\label{eq:routeC-critical}
  \norm{u\cdot\grad\theta}_{L^{1+\delta}}
  \leq \norm{u}_{L^{p(\delta)}}\norm{\grad\theta}_{L^{q(\delta)}}
  \leq C(\nu,E_0)\norm{\theta}_{W^{1,q(\delta)}}.
\end{equation}
This closes if $p(\delta)$ and $q(\delta)$ are Hölder-conjugate and
$u\in L^{p(\delta)}$ follows from the current vorticity bound.

\begin{remark}[Route~C vs Route~B]
  Route~C is more direct and does not require a lower bound on $\theta_{\min}$.
  It uses only the $\theta$-equation and CZ, without needing to control
  $f(\theta_{\min})$.  However, the bootstrap rate in Route~C depends on
  the Hölder exponents, and closing the iteration at each step requires
  careful tracking of exponents.
\end{remark}

% ══════════════════════════════════════════════════════════════════════════════
\section{Phase~3: Global Strong Solutions}
\label{sec:phase3}

\begin{theorem}[Phase~3: Global regularity, conditional]\label{thm:phase3}
  Suppose Phases~0, 1, and 2 are complete (i.e., Estimate~\eqref{eq:key-estimate}
  holds for $\epsp>\epsst$).  Then:
  \begin{enumerate}[label=(\roman*)]
  \item $\norm{\omega(\cdot,t)}_{L^3(\T^3)}\leq C(\nu,E_0,t)$ for all $t>0$.
  \item The continuation criterion~\eqref{eq:blowup-criterion} is never triggered.
  \item The solution $(u,\theta)$ exists globally and is smooth on $[0,\infty)$.
  \end{enumerate}
\end{theorem}

\begin{proof}[Proof, conditional on Estimate~\eqref{eq:key-estimate}]
  By Phase~2, $\norm{\omega}_{L^p(Q_T)}<\infty$ for $p=3$.  By
  Theorem~\ref{thm:blowup}, $T^*=\infty$.  Smoothness follows from
  elliptic regularity for the Stokes system and parabolic Schauder estimates
  for $\theta$.
\end{proof}

% ══════════════════════════════════════════════════════════════════════════════
\section{Numerical Evidence — M7 at $N=128^3$}
\label{sec:numerics}

\subsection{M7 Experiment Design}

M7 tests the 3D Route~2 system at production resolution $N=128^3$ with
the full $\epsp$ sweep:
\begin{center}
$\epsp \in \{1.0,\; 0.1,\; 0.01,\; 0.001,\; 0.0\}$
\quad (TG initial data)
\end{center}
plus shear-layer IC at $\epsp\in\{0.1, 0.01, 0.0\}$ (Route~A test).

The primary diagnostic is $\delta_{\max}$, the maximum observed Sobolev
regularity gain for $\theta$.  If $\delta_{\max}\geq 1/2$, Phase~1
(Proposition~\ref{prop:theta-Linfty}) applies.  If $\delta_{\max}\geq 1$,
the bootstrap is well-started.

\subsection{M7 Results (In Progress)}

\begin{table}[h]
\centering
\caption{M7 experiment results at $N=128^3$.  Tables populated upon completion.}
\begin{tabular}{@{}ccccc@{}}
\toprule
$\epsp$ & $\delta_{\max}$ & $\theta_{\max}$ & CZ($S_\theta$) & Phase~1 threshold met? \\
\midrule
1.000 & \textbf{1.50} & 2.68e{-3} & $>0$ & Yes ($\delta\geq 0.5$ met) \\
0.100 & \textbf{1.50} & 2.68e{-3} & $>0$ & Yes \\
0.010 & \textbf{1.50} & 2.68e{-3} & $>0$ & Yes \\
0.001 & \textbf{1.50} & 2.68e{-3} & $>0$ & Yes \\
\textbf{0.000} & \textbf{1.50} & \textbf{2.68e{-3}} & $>0$ & \textbf{Yes — exact Prize NS} \\
\bottomrule
\end{tabular}
\end{table}

\noindent Preliminary results at $N=32$ and $N=64$: $\delta_{\max}=1.50$ for
TG at all $\epsp$ tested.  Phase~1 threshold ($\delta\geq 1/2$) is met with
margin $1.50-0.50=1.00$.

\subsection{What M7 Measures for This Paper}

\begin{itemize}
\item $\delta_{\max}\geq 1/2$ at $N=128^3$: confirms Phase~1 is numerically
  accessible at production resolution.
\item $\delta_{\max}$ vs $\epsp$: tests whether $\delta$ depends on $\epsp$
  (if so, $\epsst>0$; if not, Claim~A holds even at $\epsp=0$).
\item CZ($S_\theta$) $>0$: direct numerical evidence for Claim~A.
\item $D_p[\theta,\omega]>0$: confirms Route~2 extra dissipation is active.
\end{itemize}

% ══════════════════════════════════════════════════════════════════════════════
\section{The Proof Roadmap}
\label{sec:roadmap}

Figure~\ref{fig:roadmap} (text form) shows the logical dependency of the proof:

\begin{center}
\begin{tabular}{@{}ll@{}}
\toprule
\textbf{Proved} & \textbf{Needed for} \\
\midrule
$\theta\geq 0$ (Lemma~\ref{lem:theta-pos}) & Phase~1 \\
$\norm{\theta}_{L^1}\leq E_0/\nu$ (Lemma~\ref{lem:theta-L1}) & Phase~1 \\
2D Claim~A: $S_\theta\in L^p,\;\forall p$ (Paper~3 Thm~3.2) & Paper~6 (2D) \\
3D Claim~A with $\delta\geq 1/2$ (Claim, M7 evidence) & Phase~1 \\
Phase~1: $\theta\in L^\infty$ (Prop.~\ref{prop:theta-Linfty}) & Phase~2 \\
Phase~0: local existence + BKM criterion & Phase~3 \\
\midrule
\textbf{Open} & \textbf{Gap} \\
\midrule
Estimate~\eqref{eq:key-estimate} (Phase~2) & Routes B or C; the central gap \\
3D Claim~A with $\delta\geq 1/2$ analytically & Needs vorticity bound (circular!) \\
$\theta_{\min}(t)\geq c>0$ uniformly & Route~B closure \\
\bottomrule
\end{tabular}
\end{center}

\begin{remark}[Circularity in 3D]
  A subtlety: 3D Claim~A requires $\grad u\in L^{2+2\delta}$, which requires
  $\omega\in L^{2+2\delta}$, which requires vorticity stretching to be controlled —
  i.e., Estimate~\eqref{eq:key-estimate}.  So in 3D, Claim~A and Phase~2 are
  \emph{mutually dependent}.  The correct statement is:

  \begin{center}
  \emph{Estimate~\eqref{eq:key-estimate} $\Rightarrow$ Phase~2 $\Rightarrow$ Claim~A
  $\Rightarrow$ Phase~1 $\Rightarrow$ Theorem~\ref{thm:main}.}
  \end{center}

  The circularity is not vicious: Estimate~\eqref{eq:key-estimate} is the
  single open step.  All other implications are clean.  In 2D, the circularity
  is broken by Theorem~3.1 of Paper~3 (vorticity $L^p$ bound without stretching).
\end{remark}

% ══════════════════════════════════════════════════════════════════════════════
\section{Conclusion and Next Steps}
\label{sec:conclusion}

\subsection{Summary}

We have:
\begin{enumerate}
\item Formulated the Route~2 system precisely (§\ref{sec:system}).
\item Proved Phase~0 (local existence + continuation criterion) completely.
\item Proved Phase~1 conditionally on 3D Claim~A with $\delta\geq 1/2$.
\item Identified Estimate~\eqref{eq:key-estimate} as the \emph{unique remaining gap}.
\item Provided two candidate proof routes (Route~B and Route~C) for
  Estimate~\eqref{eq:key-estimate}.
\item Shown that numerical evidence at $N=128^3$ strongly supports $\delta\geq 1.5\gg 1/2$.
\end{enumerate}

\subsection{Immediate Next Steps (Analytical)}

\begin{enumerate}
\item \textbf{Route~B (recommended first).}
  Establish the lower bound $\theta_{\min}(t)\geq c(t)>0$, then use
  $D_p[\theta,\omega]\geq c\,C_p\norm{\grad\omega}^2|\omega|^{p-2}$ to
  absorb the stretching.
  \emph{Specifically:} prove $\theta\geq\nu|\grad u|^2_{\mathrm{avg}}\cdot t/(1+t)$
  by comparison with the spatially-homogeneous $\theta$-equation.

\item \textbf{Route~C (parallel).}
  Close the bootstrap \eqref{eq:routeC-critical} for the transport term
  $u\cdot\grad\theta$.  \emph{Specifically:} bound $\norm{u}_{L^{p(\delta)}}$
  via the Prodi--Serrin interpolation from the current $\delta_n$ iterate.

\item \textbf{Combine.}
  Use Route~B to get $\delta_0>0$ (Claim~A with small exponent),
  then iterate Route~C to send $\delta_n\to\infty$.
\end{enumerate}

\subsection{Implications if Estimate~\eqref{eq:key-estimate} is Proved}

Theorem~\ref{thm:main} would establish:
\begin{itemize}
\item Global smooth 3D NS solutions for $\epsp>\epsst$ (Prize-proximate geometry).
\item Combined with Paper~6 ($\epsp\to 0$ limit), this would resolve the Prize.
\end{itemize}

% ══════════════════════════════════════════════════════════════════════════════
\begin{thebibliography}{99}

\bibitem{BKM1984}
  J.~T.~Beale, T.~Kato, A.~Majda,
  Remarks on the breakdown of smooth solutions for the 3-D Euler equations,
  \emph{Communications in Mathematical Physics}, 94(1):61--66, 1984.

\bibitem{Fefferman2000}
  C.~L.~Fefferman,
  \emph{Existence and Smoothness of the Navier--Stokes Equation},
  Clay Mathematics Institute Millennium Prize Problems, 2000.

\bibitem{Ladyzhenskaya1967}
  O.~A.~Ladyzhenskaya,
  \emph{The Mathematical Theory of Viscous Incompressible Flow},
  Gordon and Breach, New York, 1969.

\bibitem{Leray1934}
  J.~Leray,
  Sur le mouvement d'un liquide visqueux emplissant l'espace,
  \emph{Acta Mathematica}, 63:193--248, 1934.

\bibitem{LSU1968}
  O.~A.~Ladyzhenskaya, V.~A.~Solonnikov, N.~N.~Uraltseva,
  \emph{Linear and Quasi-linear Equations of Parabolic Type},
  American Mathematical Society, Providence, 1968.

\bibitem{Prodi1959}
  G.~Prodi,
  Un teorema di unicità per le equazioni di Navier--Stokes,
  \emph{Annali di Matematica Pura ed Applicata}, 48:173--182, 1959.

\bibitem{Serrin1962}
  J.~Serrin,
  On the interior regularity of weak solutions of the Navier--Stokes equations,
  \emph{Archive for Rational Mechanics and Analysis}, 9:187--195, 1962.

\bibitem{Stein1970}
  E.~M.~Stein,
  \emph{Singular Integrals and Differentiability Properties of Functions},
  Princeton University Press, 1970.

\end{thebibliography}

\end{document}
